\section{The Coherent Side for Toric Varieties}

In this section we specialize the coherent side of the
coherent--constructible correspondence to the toric setting.  Let
\(X_\Sigma\) be a smooth toric variety with acting torus \(T\), and let
\(M\) and \(N\) denote the character and cocharacter lattices of \(T\).
The coherent side of the correspondence is modeled by the dg category
\[
\operatorname{Perf}_T(X_\Sigma),
\]
or equivalently, in the smooth setting, by the bounded equivariant
coherent derived category \(D^b(\operatorname{Coh}_T(X_\Sigma))\).

The toric structure gives a particularly explicit description of this
category.  On an affine toric chart
\(U_\sigma=\operatorname{Spec}k[\sigma^\vee\cap M]\), \(T\)-equivariant
quasi-coherent sheaves correspond to \(M\)-graded modules over the
semigroup algebra \(k[\sigma^\vee\cap M]\).  Thus equivariance translates
the torus action into a grading by the character lattice \(M\).

The purpose of this section is to recall this graded description and to
introduce the coherent generators used in the proof of the toric CCC.
These generators are indexed by pairs \((\sigma,\chi)\), where
\(\sigma\in\Sigma\) is a cone of the fan and \(\chi\in M\) is a
character.  They will later be matched with standard constructible
sheaves on \(M_{\mathbb R}\).  The main result prepared in this section
is that these coherent objects generate
\(\operatorname{Perf}_T(X_\Sigma)\).


\subsection{Equivariant Coherent Sheaves on Toric Varieties}

Let \(\sigma\in\Sigma\) be a cone, and let
\(U_\sigma\subset X_\Sigma\) be the corresponding affine toric open
subset.  We write \(U_\sigma=\operatorname{Spec}A_\sigma\), where
\(A_\sigma=k[\sigma^\vee\cap M]\).  The semigroup algebra \(A_\sigma\)
has its natural \(M\)-grading
\[
A_\sigma=\bigoplus_{m\in\sigma^\vee\cap M}A_{\sigma,m},
\qquad
A_{\sigma,m}=k\cdot\chi^m.
\]
Equivalently, \(\deg(\chi^m)=m\).  This grading records the torus action
on \(U_\sigma\).

We denote by \(A_\sigma\text{-}\operatorname{Mod}^M\) the category of
\(M\)-graded \(A_\sigma\)-modules.  Thus an object is an
\(A_\sigma\)-module \(P=\bigoplus_{\chi\in M}P_\chi\) such that
\(A_{\sigma,m}\cdot P_\chi\subset P_{\chi+m}\) for all
\(m\in\sigma^\vee\cap M\) and \(\chi\in M\).

\begin{proposition}\label{prop:affine-equiv-graded}
There is an equivalence of categories
\[
\operatorname{QCoh}_T(U_\sigma)
\simeq
A_\sigma\text{-}\operatorname{Mod}^{M}.
\]
\end{proposition}

\begin{proof}
Since \(U_\sigma\) is affine, a quasi-coherent sheaf is determined by its
module of global sections.  If \(\mathcal F\) is \(T\)-equivariant, then
\(\Gamma(U_\sigma,\mathcal F)\) carries a compatible \(T\)-action.  Since
\(T\) is a torus, this module decomposes into weight spaces:
\[
\Gamma(U_\sigma,\mathcal F)
=
\bigoplus_{\chi\in M}\Gamma(U_\sigma,\mathcal F)_\chi .
\]
Compatibility with the \(A_\sigma\)-module structure is precisely the
grading condition
\[
A_{\sigma,m}\cdot \Gamma(U_\sigma,\mathcal F)_\chi
\subset
\Gamma(U_\sigma,\mathcal F)_{\chi+m}.
\]
Thus \(\Gamma(U_\sigma,\mathcal F)\) is an \(M\)-graded
\(A_\sigma\)-module.

Conversely, an \(M\)-graded \(A_\sigma\)-module \(P\) defines a
quasi-coherent sheaf \(\widetilde P\) on \(U_\sigma\), and the grading
defines a \(T\)-equivariant structure on \(\widetilde P\).  These two
constructions are inverse to each other.
\end{proof}

Let \(A_\sigma\text{-}\operatorname{mod}^M\) denote the full subcategory
of finitely generated \(M\)-graded \(A_\sigma\)-modules.

\begin{corollary}
There is an equivalence of categories
\[
\operatorname{Coh}_T(U_\sigma)
\simeq
A_\sigma\text{-}\operatorname{mod}^{M}.
\]
\end{corollary}

\begin{proof}
By \cref{prop:affine-equiv-graded}, \(T\)-equivariant
quasi-coherent sheaves on \(U_\sigma\) are equivalent to \(M\)-graded
\(A_\sigma\)-modules.  Since \(U_\sigma\) is affine, coherence is
equivalent to finite generation over \(A_\sigma\).  Hence the coherent
objects correspond exactly to finitely generated graded modules.
\end{proof}



We now pass from a single affine chart to the whole toric variety.  The
fan \(\Sigma\) gives an affine open cover
\(X_\Sigma=\bigcup_{\sigma\in\Sigma}U_\sigma\), and for two cones
\(\sigma,\tau\in\Sigma\), one has
\(U_\sigma\cap U_\tau=U_{\sigma\cap\tau}\).  The open immersions
\(U_{\sigma\cap\tau}\hookrightarrow U_\sigma\) and
\(U_{\sigma\cap\tau}\hookrightarrow U_\tau\) induce morphisms of
\(M\)-graded algebras \(A_\sigma\to A_{\sigma\cap\tau}\) and
\(A_\tau\to A_{\sigma\cap\tau}\).




\begin{proposition}[Gluing over the fan]
A \(T\)-equivariant quasi-coherent sheaf on \(X_\Sigma\) is equivalently
given by a family of \(M\)-graded \(A_\sigma\)-modules \(P_\sigma\), one
for each cone \(\sigma\in\Sigma\), together with isomorphisms of
\(M\)-graded \(A_{\sigma\cap\tau}\)-modules
\[
A_{\sigma\cap\tau}\otimes_{A_\sigma}P_\sigma
\xrightarrow{\sim}
A_{\sigma\cap\tau}\otimes_{A_\tau}P_\tau
\]
for all pairs \(\sigma,\tau\in\Sigma\), satisfying the usual cocycle
condition on triple overlaps.  The corresponding sheaf is coherent if
and only if each \(P_\sigma\) is a finitely generated \(M\)-graded
\(A_\sigma\)-module.
\end{proposition}

\begin{proof}
This is the usual descent description of quasi-coherent sheaves with
respect to the affine open cover \(\{U_\sigma\}_{\sigma\in\Sigma}\),
combined with the equivalence
\(\operatorname{QCoh}_T(U_\sigma)\simeq
A_\sigma\text{-}\operatorname{Mod}^M\) on each affine chart.  Restricting
a sheaf from \(U_\sigma\) to \(U_{\sigma\cap\tau}\) corresponds
algebraically to extension of scalars along
\(A_\sigma\to A_{\sigma\cap\tau}\).  Since these restriction maps
preserve the \(M\)-grading, the gluing maps are isomorphisms of graded
modules.  The cocycle condition is exactly the usual sheaf gluing
condition on triple overlaps.  Finally, coherence is local on
\(X_\Sigma\), so the coherent case is obtained by requiring each
\(P_\sigma\) to be finitely generated.
\end{proof}

Thus \(T\)-equivariant coherent sheaves on \(X_\Sigma\) can be described
in terms of \(M\)-graded modules over the affine semigroup algebras
\(A_\sigma\), together with compatibility over the fan.  This
local-to-global description will be used in the construction of the
coherent generators.



\subsection{Coherent Generators}

We now introduce the coherent generators used in the toric
coherent--constructible correspondence.  For each cone
\(\sigma\in\Sigma\), let \(i_\sigma:U_\sigma\hookrightarrow X_\Sigma\)
be the inclusion of the corresponding affine toric open subset. 


\begin{definition}
Let \(\sigma\in\Sigma\), and let \(U_\sigma\) be the corresponding
affine toric open subset.  For a character \(\chi\in M\), define
$
\mathcal O_{U_\sigma}(\chi)
$
to be the trivial line bundle \(U_\sigma\times k\to U_\sigma\) with the
following \(T\)-action on its total space:
\[
t\cdot(x,v)=(tx,\chi(t)v),
\qquad
t\in T,\ x\in U_\sigma,\ v\in k.
\]
\end{definition}

\begin{proposition}
The line bundle \(\mathcal O_{U_\sigma}(\chi)\) is a
\(T\)-equivariant line bundle on \(U_\sigma\).
\end{proposition}

\begin{proof}
Since \(\chi:T\to k^*\) is a character, it is a group homomorphism.
Hence, for \(t_1,t_2\in T\), one has
\[
t_1\cdot(t_2\cdot(x,v))
=
t_1\cdot(t_2x,\chi(t_2)v)
=
(t_1t_2x,\chi(t_1)\chi(t_2)v)
=
(t_1t_2x,\chi(t_1t_2)v).
\]
Thus the formula defines a \(T\)-action on \(U_\sigma\times k\).

Let \(\pi:U_\sigma\times k\to U_\sigma\) be the projection.  Then
\[
\pi(t\cdot(x,v))=\pi(tx,\chi(t)v)=tx=t\cdot\pi(x,v),
\]
so \(\pi\) is \(T\)-equivariant.  Moreover, for each \(x\in U_\sigma\),
the induced map on fibers
\[
(\mathcal O_{U_\sigma}(\chi))_x
\longrightarrow
(\mathcal O_{U_\sigma}(\chi))_{tx}
\]
is multiplication by \(\chi(t)\), hence is \(k\)-linear.  Therefore
\(\mathcal O_{U_\sigma}(\chi)\) is a \(T\)-equivariant line bundle.
\end{proof}

We call this equivariant structure the \textbi{\(T\)-linearization of weight
\(\chi\)}.  Thus \(\mathcal O_{U_\sigma}(\chi)\) has the same underlying
line bundle as \(\mathcal O_{U_\sigma}\), but its \(T\)-equivariant
structure is twisted by the character \(\chi\).

\begin{remark}
The word ``linearization'' means that the \(T\)-action on
\(U_\sigma\) is lifted to the total space of the line bundle, with linear
maps on fibers.  Here the induced map
\[
(\mathcal O_{U_\sigma})_x\longrightarrow
(\mathcal O_{U_\sigma})_{tx}
\]
is multiplication by \(\chi(t)\).  Equivalently, this is descent data for
the action groupoid \(T\times U_\sigma\rightrightarrows U_\sigma\).
\end{remark}


We define
\[
\Theta'(\sigma,\chi)
:=
(i_\sigma)_*\mathcal O_{U_\sigma}(\chi).
\]
These objects are \(T\)-equivariant quasi-coherent sheaves on \(X_\Sigma\).
They are indexed by the same combinatorial data \((\sigma,\chi)\) which
will later index the constructible generators.




\begin{theorem}[Theta generation in the equivariant quasi-coherent category]\label{thm:coherent-theta-generation}
Let \(X_\Sigma\) be a smooth toric variety. Then
\[
D(\operatorname{QCoh}_T(X_\Sigma))
=
\Big\langle
\Theta'(\sigma,\chi)
\ \Big|\
\sigma\in\Sigma,\ \chi\in M
\Big\rangle_{\mathrm{loc}}.
\]
Here \(\langle - \rangle_{\mathrm{loc}}\) denotes the smallest localizing
triangulated subcategory containing the indicated objects.
\end{theorem}

\begin{proof}
Let \(\mathcal C\) be the right-hand side. We prove that every object of
\(D(\operatorname{QCoh}_T(X_\Sigma))\) belongs to \(\mathcal C\).

First consider an affine toric chart
\(U_\sigma=\operatorname{Spec}A_\sigma\), where
\(A_\sigma=k[\sigma^\vee\cap M]\). By \cref{prop:affine-equiv-graded}, one has
\(\operatorname{QCoh}_T(U_\sigma)\simeq A_\sigma\operatorname{-Mod}^M\).
Under this equivalence, \(\mathcal O_{U_\sigma}(\chi)\) corresponds to the
free graded module \(A_\sigma(\chi)\). The modules \(A_\sigma(\chi)\), for
\(\chi\in M\), generate \(D(A_\sigma\operatorname{-Mod}^M)\): if
\(P\in D(A_\sigma\operatorname{-Mod}^M)\) satisfies
\(\operatorname{RHom}_{A_\sigma}(A_\sigma(\chi),P)=0\) for every
\(\chi\in M\), then all graded components of all cohomology modules
\(H^i(P)\) vanish, hence \(P\simeq 0\).

Now let \(F\in D(\operatorname{QCoh}_T(X_\Sigma))\). Choose the finite
\(T\)-stable affine cover of \(X_\Sigma\) by affine toric charts. Since
intersections of affine toric charts are again affine toric charts, the
associated \v{C}ech resolution expresses \(F\) as a finite totalization of
objects of the form \((i_\tau)_*(F|_{U_\tau})\), with \(\tau\in\Sigma\).
For each \(\tau\), the restriction \(F|_{U_\tau}\) is generated by
\(\mathcal O_{U_\tau}(\chi)\), \(\chi\in M\). Since \((i_\tau)_*\) is exact
and preserves coproducts for these quasi-compact open immersions,
\((i_\tau)_*(F|_{U_\tau})\) is generated by
\((i_\tau)_*\mathcal O_{U_\tau}(\chi)=\Theta'(\tau,\chi)\). Hence every term
in the \v{C}ech resolution belongs to \(\mathcal C\).

Since \(\mathcal C\) is localizing, it is closed under shifts, cones, finite
totalizations, and coproducts. Therefore \(F\in\mathcal C\). This proves
\[
D(\operatorname{QCoh}_T(X_\Sigma))
=
\big\langle \Theta'(\sigma,\chi)\mid \sigma\in\Sigma,\ \chi\in M
\big\rangle_{\mathrm{loc}}.
\]
\end{proof}

\begin{corollary}[Compact form of the coherent side]
\label{thm:compact-form-of-the-coherent-side}
Let \(X_\Sigma\) be a smooth toric variety. Then the compact part of the
equivariant quasi-coherent derived category is the perfect equivariant category:
\[
D(\operatorname{QCoh}_T(X_\Sigma))^c
\simeq
\operatorname{Perf}_T(X_\Sigma).
\]
\end{corollary}

\begin{proof}
This follows directly from
\cref{thm:coherent-theta-generation} and
\cref{thm:perfect-complexes-as-compact-objects}.
\end{proof}

Thus the equivariant quasi-coherent side has a distinguished family of theta
generators indexed by pairs \((\sigma,\chi)\). The perfect coherent side will be
obtained by passing to compact objects as in
\cref{thm:compact-form-of-the-coherent-side}. 